% Figura Anexo I: el diagrama original de HICKS (1937), adaptado (curvas IS y LL, notación de Hicks: renta I, interés i)
\documentclass[tikz,border=3pt]{standalone}
\usepackage[T1]{fontenc}
\usepackage[utf8]{inputenc}
\usetikzlibrary{arrows.meta,calc}
\begin{document}
\begin{tikzpicture}[>={Latex[length=2.2mm]}, line cap=round, font=\footnotesize,
  eje/.style={->, thick},
  curva/.style={very thick},
  nueva/.style={very thick, blue!55!black},
  guia/.style={densely dashed, gray!80},
  pto/.style={circle, fill=black, inner sep=1.1pt}]
  % Ejes
  \draw[eje] (0,0) -- (10.6,0) node[below left=1pt]{Renta, $I$};
  \draw[eje] (0,0) -- (0,6.1) node[below right=2pt]{Tipo de interés, $i$};
  % Curva LL: tramo plano (trampa de la liquidez) y tramo creciente
  \draw[curva] (0.5,0.95) -- (2.9,0.95) to[out=4,in=215] (5,2.1) to[out=35,in=225] (7.8,4.4) to[out=45,in=250] (8.9,5.5) node[right]{$LL$};
  % Curva IS principal
  \draw[curva] (1.3,5.4) to[out=-52,in=160] (5,2.1) to[out=-20,in=172] (9.3,1.15) node[right]{$IS$};
  % IS deprimida que corta el tramo plano
  \draw[nueva, densely dashed] (0.7,3.5) to[out=-62,in=140] (1.9,0.95) to[out=-40,in=130] (2.9,0.15) node[below right=-2pt]{$IS'$};
  % Equilibrio P
  \node[pto] (P) at (5,2.1) {};
  \node[above right=0pt of P, inner sep=1pt] {$P$};
  \draw[guia] (5,2.1) -- (5,0) node[below]{$I_P$};
  \draw[guia] (5,2.1) -- (0,2.1) node[left]{$i_P$};
  % Equilibrio deprimido P'
  \node[pto, fill=blue!55!black] (Pp) at (1.9,0.95) {};
  \node[above=1pt of Pp, inner sep=1pt, blue!55!black] {$P'$};
  % Anotaciones
  \draw[<-, gray!60!black] (1.6,0.8) to[out=-60,in=175] (3.1,-0.85) node[right, align=left, text=black]{tramo plano: preferencia absoluta por la liquidez\\ (atesoramiento); el interés no puede bajar más};
  \draw[<-, gray!60!black] (8.2,4.6) to[out=160,in=-30] (6.6,5.5) node[above left=-3pt, align=right, text=black]{tramo creciente: la renta eleva la demanda\\ de dinero y, con $M$ dada, el interés};
\end{tikzpicture}
\end{document}
